% ==========================================================================
% resumoPt.tex
% Resumo em língua vernácula (NBR 6028)
% ==========================================================================

\begin{center}
\ABNTEXchapterfont\bfseries\Large\imprimirtitulo
\end{center}

\begin{resumo}
A democratização do acesso à baixa órbita terrestre por meio de nanossatélites da classe CubeSat impõe desafios rigorosos de engenharia de sistemas. Durante o lançamento e a ejeção, o chassi estrutural é submetido a carregamentos dinâmicos severos, como o choque transiente na interface do \textit{Poly-Picosatellite Orbital Deployer} (P-POD) e a vibração aleatória acústica regulamentada pela norma NASA GSFC-STD-7000A (GEVS). Chassis monolíticos convencionais apresentam elevada massa e transferem altos níveis de aceleração à carga útil devido ao acoplamento dinâmico direto. Este trabalho propõe uma rota computacional de alto desempenho para a concepção, a modelagem paramétrica e a otimização multiobjetivo de um chassi CubeSat 1U com painéis constituídos por metamateriais auxéticos reentrantes em liga aeroespacial AlSi10Mg, fabricado por Fusão em Leito de Pó a Laser (L-PBF). Um modelo paramétrico tridimensional, sujeito a critérios estritos de Manufatura Aditiva (DfAM), foi integrado a um modelo elastodinâmico multiescala calibrado por elementos finitos para gerar uma campanha de 250 pontos amostrais via Hipercubo Latino. Treinou-se um modelo substituto neural residual guiado por física (\textit{Physics-Guided ResNet}) operando sobre vetor de saída desacoplado em $\mathbb{R}^3$ ($\hat{f}_1, \hat{\sigma}_{3\sigma}, \hat{G}_{\text{rms, payload}}$) e pós-processamento determinístico da transmissibilidade ($\hat{T} \equiv \hat{G}_{\text{rms, payload}} / \num{14.1}$), eliminando redundâncias lineares. A rede, regularizada por perdas de Gibson-Ashby e monotonicidade de rigidez, alcançou $\overline{R^2} = \num{0.9972} \pm \num{0.0016}$ e $\text{NRMSE} = \SI{1.27 \pm 0.34}{\percent}$ em validação cruzada estratificada $10$-fold sobre as 250 amostras, com latência de inferência de \SI{0.405}{\milli\second} (\textit{speedup} instantâneo superior a \num{200000} vezes e \textit{speedup} amortizado de ponta a ponta de $\approx \num{40.0}\times$). Uma máquina de estados finitos determinística, implementada em LangGraph como grafo de fluxo de dados fortemente tipado, orquestrou o descarte precoce de candidatos inviáveis. A otimização multiobjetivo pelo algoritmo NSGA-II, com \num{10000} avaliações de função, mapeou a fronteira de Pareto entre massa e transmissibilidade. Pelo método multicritério TOPSIS, elegeu-se o projeto ótimo de voo (\#28: $\theta^* = \SI{65.75}{\degree}$ com deflexão cartesiana $\theta_{\text{def}} = \SI{-21.40}{\degree}$, $t^* = \SI{0.613}{\milli\metre}$, cotas entre eixos nodais $l^* = \SI{10.42}{\milli\metre}$ e $h^* = \SI{14.85}{\milli\metre}$, comprimento livre da costela $l_{\text{free}} = \SI{6.42}{\milli\metre}$, Poisson micromecânico 2D $\nu_{yx}^* = \num{-12.63}$ e Poisson estrutural 3D homogeneizado $\nu_{\text{eff}} = \num{-0.428}$). A validação numérica de alta ordem por elementos finitos tridimensionais (\textit{Code\_Aster Ground Truth}) confirmou erro de apenas \SI{0.99}{\percent} na tensão de pico e de \SI{0.008}{\percent} na transmissibilidade predita. Em relação ao chassi monolítico convencional, a estrutura auxética otimizada proporcionou redução de massa de \SI{62.49}{\percent} (\SI{0.308}{\kilo\gram} $\to$ \SI{0.116}{\kilo\gram}), atenuação elastodinâmica passiva em regime sub-comprimento de onda de \SI{76.50}{\percent} na transmissibilidade vibratória (\SI{12.6}{\decibel} por antirressonância mecânica e descasamento de impedância acústica $> \SI{66}{\percent}$, invariante a variações de amortecimento modal $\zeta \in \{1{,}0\%, 2{,}0\%, 3{,}5\%\}$), frequência fundamental de \SI{579.62}{\hertz} e pleno atendimento aos critérios de escoamento e de fadiga estocástica de Steinberg no entalhe crítico ($D = \num{0.0458} \ll \num{0.25}$), com integridade subcrítica chancelada pelo critério de Murakami ($\Delta K = \SI{1.28}{\mega\pascal\sqrt{\metre}} < \Delta K_{th} = \SI{2.50}{\mega\pascal\sqrt{\metre}}$) sob rotas de alívio térmico a $\SI{300}{\celsius}$ e $\SI{250}{\celsius}$, atingindo prontidão Digital TRL 4 (NASA SP-6105) e TRL 3 de hardware (ISO 16290:2013).

\vspace{\onelineskip}
\noindent
\textbf{Palavras-chave}: CubeSat 1U. Metamateriais auxéticos. \textit{Design for Additive Manufacturing}. Redes neurais guiadas por física. Otimização multiobjetivo NSGA-II. NASA GEVS.
\end{resumo}
